\documentclass[11pt]{article}
\usepackage{amsmath, amssymb}
\title{Audited Mathematical Notes on Partition Functions and SAT Complexity}
\date{2026-11-24}
\begin{document}
\maketitle

\section{Definitions}
For a 3-CNF formula $\varphi$ over $n$ variables with $m$ clauses, let $E_\varphi(y) \in \{0,\dots,m\}$ denote the number of clauses violated by assignment $y \in \{0,1\}^n$.

\subsection{Partition Function}
\begin{equation}
Z_\varphi(z) = \sum_{y \in \{0,1\}^n} z^{E_\varphi(y)} = \sum_{k=0}^{m} c_k z^k, \qquad c_k = \#\{y \in \{0,1\}^n : E_\varphi(y) = k\}.
\end{equation}
Note: $Z_\varphi(0) = c_0 = \#\mathrm{SAT}(\varphi)$, which is $\#P$-complete to compute. Furthermore, $Z_\varphi(z)$ is not invariant under logically equivalent syntactic transformations of $\varphi$ (e.g., duplicating a clause shifts $E_\varphi(y)$).

\subsection{Quantum / Trigonometric Amplitude}
\begin{equation}
A_\varphi(\theta) = 2^{-n} Z_\varphi(e^{i\theta}) = \mathbb{E}_{y \sim U}\!\left[e^{i\theta E_\varphi(y)}\right].
\end{equation}
At $\theta = \pi$, $A_\varphi(\pi) = 2^{-n}(N_{\mathrm{even}} - N_{\mathrm{odd}})$. Monte Carlo estimation over $N_s$ uniform samples incurs statistical noise $\Theta(1/\sqrt{N_s})$, which overwhelms signals of magnitude $2^{-\Omega(n)}$ unless $N_s = \Omega(2^{2n})$.

\subsection{Complex Zero Radius and Inverse Temperature}
With inverse temperature $\beta \in \mathbb{C}$, $Z_\varphi(\beta) = \sum_{x \in \{0,1\}^n} e^{-\beta E_\varphi(x)}$ (setting $z = e^{-\beta}$). The Lee--Yang zero radius is
\begin{equation}
d(\varphi) = \min \{ |z| : Z_\varphi(z) = 0 \}, \qquad Z_\varphi(i) = \sum_{x \in \{0,1\}^n} i^{E_\varphi(x)}.
\end{equation}
Note: $Z_\varphi(z)$ is syntactically non-invariant: $\varphi = (x)$ yields $Z_\varphi(z) = 1+z$ ($z = -1$), whereas the logically equivalent $\psi = (x) \land (x)$ yields $Z_\psi(z) = 1+z^2$ ($z = \pm i$).

\subsection{Galois Group of the Partition Polynomial}
Let $G_\varphi = \mathrm{Gal}(Z_\varphi / \mathbb{Q})$. Because $Z_\varphi(z) = \sum_{k=0}^m c_k z^k$, $G_\varphi$ depends solely on the degeneracy counts $(c_0, \dots, c_m)$, not on $\mathrm{Aut}(\varphi)$ in $\{0,1\}^n$. Generic integer polynomials have $G_\varphi \cong S_m$, while $Z(z) = 1 + z^m$ yields the abelian cyclotomic group $(\mathbb{Z}/2m\mathbb{Z})^\times$.

\section{Non-uniform Advice Bounds (Karp--Lipton)}
If $\mathrm{NP} \subseteq \mathrm{P}/\mathrm{poly}$ (polynomial-size advice/strategy guide), then the polynomial hierarchy collapses to its second level:
\begin{equation}
\mathrm{NP} \subseteq \mathrm{P}/\mathrm{poly} \implies \Sigma_2^P = \Pi_2^P.
\end{equation}
This does not imply $\mathrm{P} = \mathrm{NP}$, and fixed $O(\log n)$-bit advice cannot specify general $n$-bit satisfying assignments.

\section{Entropy and Algebraic Complexity (VP vs VNP)}
\subsection{Violation-Count Shannon Entropy}
With $p_k = c_k / 2^n$, the Shannon entropy of the violation count distribution is
\begin{equation}
H(E_\varphi) = -\sum_{k=0}^{m} p_k \log_2 p_k.
\end{equation}
Note: Computing $p_k$ exactly requires $\#P$-hard counting of $c_k$.

\subsection{Multilinear Polynomial Family in VNP}
Viewing clause structure as formal variables $X$, the induced multilinear polynomial family $\{F_n(X, z)\}_{n \ge 1}$ satisfies
\begin{equation}
\{F_n\}_{n \ge 1} \in \mathrm{VNP}.
\end{equation}
Note: Membership $\{F_n\} \in \mathrm{VNP}$ follows directly from the definition of $\mathrm{VNP}$ and does not establish $F_n \notin \mathrm{VP}$. Moreover, a separation $\mathrm{VP} \neq \mathrm{VNP}$ over a field of characteristic $0$ does not directly imply $\mathrm{P} \neq \mathrm{NP}$ for Boolean circuits.

\section{Information Theory and Proof Complexity}
\subsection{Search Surprisal and Mutual Information with Advice}
For $K = \#\mathrm{SAT}(\varphi)$ satisfying assignments, the uniform winning probability and surprisal are
\begin{equation}
p_{\mathrm{win}} = \frac{K}{2^n}, \qquad I_{\mathrm{win}} = -\log_2 p_{\mathrm{win}} = n - \log_2 K, \qquad \mathbb{E}[N_{\mathrm{trials}}] = \frac{2^n}{K} = 2^{I_{\mathrm{win}}}.
\end{equation}
With advice/debug history $D$ and instance $X$, the conditional mutual information is
\begin{equation}
I(Y; D \mid X) = H(Y \mid X) - H(Y \mid X, D).
\end{equation}

\subsection{Propositional Proof Length}
For a proof system $P$ and tautology $\tau$, the shortest proof length is $L_P(\tau) = \min_{\pi : P(\pi)=\tau} |\pi|$. By Cook--Reckhow, $\mathrm{NP} = \mathrm{coNP}$ holds if and only if there exists a polynomially bounded propositional proof system.

\end{document}
